\documentclass[journal, onecolumn]{IEEEtran}

% --- Codificacion y lenguaje -------------------------------------------------
\usepackage[utf8]{inputenc}
\usepackage[T1]{fontenc}
\usepackage[spanish, es-tabla]{babel}

% --- Margenes ----------------------------------------------------------------
\usepackage[top=2.5cm, bottom=2.5cm, left=2.5cm, right=2.5cm]{geometry}

\usepackage{listings}
\usepackage{xcolor}

\lstdefinestyle{ArduinoStyle}{
    language=C++,
    basicstyle=\ttfamily\small,
    keywordstyle=\color{blue},
    commentstyle=\color{green!60!black},
    stringstyle=\color{orange},
    numbers=left,
    numberstyle=\tiny,
    frame=single,
    breaklines=true,
    backgroundcolor=\color{gray!10}
}

% --- Graficos y figuras ------------------------------------------------------
\usepackage{graphicx}
\usepackage{caption}
\usepackage{float}

% --- Colores -----------------------------------------------------------------
\usepackage{xcolor}
\definecolor{azulTitulo}{RGB}{0,70,127}
\definecolor{codegreen}{rgb}{0,0.6,0}
\definecolor{codegray}{rgb}{0.5,0.5,0.5}
\definecolor{codepurple}{rgb}{0.58,0,0.82}
\definecolor{backcolour}{rgb}{0.95,0.95,0.92}

% --- Tablas ------------------------------------------------------------------
\usepackage{booktabs}
\usepackage{tabularx}
\usepackage{array}

% --- Hipervinculos -----------------------------------------------------------
\usepackage[hidelinks]{hyperref}

% --- Encabezado y pie de pagina ----------------------------------------------
\usepackage{fancyhdr}
\setlength{\headheight}{33pt}
\pagestyle{fancy}
\fancyhf{}
\lhead{%
  \begin{picture}(0,0)
    \put(5,10){\includegraphics[width=11mm]{pictures/escudo.png}}
  \end{picture}%
  \hspace{1mm}\textsc{\color{azulTitulo}Control Digital \hspace{90mm} Ing. Telecomunicaciones}%
}
\rhead{%
  \begin{picture}(125,0)
    \put(10,10){\includegraphics[width=35mm]{pictures/logo.png}}
  \end{picture}%
}
\cfoot{\thepage}
\renewcommand{\headrulewidth}{0.6pt}
\renewcommand{\sectionmark}[1]{}
\renewcommand{\subsectionmark}[1]{}

\fancypagestyle{plain}{
  \fancyhf{}
  \lhead{%
    \begin{picture}(0,0)
      \put(0,-5){\includegraphics[width=8.5mm]{pictures/escudo.png}}
    \end{picture}%
  }
  \rhead{%
    \begin{picture}(86,0)
      \put(0,-10){\includegraphics[width=30mm]{pictures/logo.png}}
    \end{picture}%
  }
  \cfoot{\thepage}
  \renewcommand{\headrulewidth}{0.6pt}
}

% --- Formato de secciones ----------------------------------------------------
\usepackage{titlesec}
\usepackage{parskip}

\titleformat{\section}
  {\color{azulTitulo}\large\bfseries}{\thesection.}{0.5em}{}
  [\color{azulTitulo}\titlerule]

\titleformat{\subsection}
  {\color{azulTitulo}\normalsize\bfseries}{\thesubsection}{0.5em}{}[]

\titleformat{\subsubsection}
  {\color{azulTitulo}\normalsize\itshape}{\thesubsubsection}{0.5em}{}[]

% --- Numeracion de secciones -------------------------------------------------
\renewcommand{\thesection}{\arabic{section}}
\renewcommand{\thesubsection}{\thesection.\arabic{subsection}}
\renewcommand{\thesubsubsection}{\thesubsection.\arabic{subsubsection}}

% --- Codigo fuente -----------------------------------------------------------
\usepackage{listings}
\lstdefinestyle{mystyle}{
    backgroundcolor=\color{backcolour},
    commentstyle=\color{codegreen},
    keywordstyle=\color{magenta},
    numberstyle=\tiny\color{codegray},
    stringstyle=\color{codepurple},
    basicstyle=\ttfamily\footnotesize,
    breakatwhitespace=false,
    breaklines=true,
    captionpos=b,
    keepspaces=true,
    numbers=left,
    numbersep=5pt,
    showspaces=false,
    showstringspaces=false,
    showtabs=false,
    tabsize=2
}
\lstset{style=mystyle}

% --- Diagramas TikZ ----------------------------------------------------------
\usepackage{tikz}
\usetikzlibrary{arrows.meta, shapes.geometric, positioning, fit, backgrounds, babel}

% --- Matematicas -------------------------------------------------------------
\usepackage{amsmath}
\usepackage{amssymb}
\usepackage{enumerate}

% --- Nombres localizados -----------------------------------------------------
\renewcommand{\figurename}{Fig.}
\renewcommand{\tablename}{Tabla}
\renewcommand{\IEEEkeywordsname}{Palabras Clave}
\addto\captionsspanish{\renewcommand{\appendixname}{Anexo}}
\newcolumntype{C}[1]{>{\centering\arraybackslash}m{#1}}

% =============================================================================
%  PORTADA
% =============================================================================

\title{%
  \color{azulTitulo}\textbf{Practica 3: Control de Motor DC}\\[4pt]
}
\author{%
  \begin{tabular}{c c}
    Franklin Benalcazar S. & Luis Felipe Rivas P.\\
    \color{azulTitulo}\underline{franklin.benalcazar@ucuenca.edu.ec} &
    \color{azulTitulo}\underline{luisf.rivas@ucuenca.edu.ec} \\
  \end{tabular}\\
  \textbf{Ingenieria en Telecomunicaciones}\\
  \textbf{Universidad de Cuenca}\\[4pt]
}

\begin{document}
\maketitle

% --- Resumen -----------------------------------------------------------------
\begin{abstract}
    La práctica consistió en diseñar e implementar un sistema de control de velocidad para un motor DC utilizando señales PWM y retroalimentación mediante un encoder óptico de 36 pulsos por revolución. Se configuró el hardware del EPC y el Controlino Mega para adquirir los pulsos del encoder y procesar la información mediante interrupciones externas y temporizadores en modo CTC.

    Posteriormente, se desarrolló una interfaz HMI para la visualización en tiempo real de variables como velocidad y referencia, permitiendo además el ajuste de parámetros de control. Se implementó la comunicación entre el controlador y la interfaz, logrando el intercambio de datos necesario para el funcionamiento del sistema PID.

    Finalmente, se verificó el comportamiento del motor ante diferentes referencias de velocidad, observando la respuesta dinámica del sistema y evaluando la precisión del control implementado.
\end{abstract}

\begin{IEEEkeywords}
Controllino Mega, HMI Stone, PWM, PID, comunicacion serial, retardo no bloqueante, control digital
\end{IEEEkeywords}

% =============================================================================
\section{Introduccion}
% =============================================================================

En esta práctica de laboratorio se desarrolló un sistema de monitoreo y control de velocidad para un motor de corriente continua (DC) utilizando el Controlino Mega y una interfaz HMI. El trabajo se centró en la adquisición de señales provenientes de un encoder óptico, el procesamiento de datos mediante interrupciones y la generación de señales PWM para regular la velocidad del motor. Además, se implementó un sistema de comunicación entre el controlador y la interfaz gráfica para visualizar variables del proceso y ajustar parámetros de control.

La práctica permitió aplicar conceptos fundamentales de control digital, temporización, adquisición de datos y programación embebida. También se analizaron aspectos relacionados con el uso de interrupciones externas y temporizadores para garantizar un muestreo preciso de la velocidad del motor, así como la integración de un controlador PID para mejorar el desempeño dinámico del sistema.

% =============================================================================
\section{Marco Teorico}
% =============================================================================

\subsection{Controladores PID}

El controlador Proporcional–Integral–Derivativo (PID) es uno de los algoritmos de control más utilizados en la industria debido a su simplicidad, robustez y capacidad de adaptarse a una gran variedad de procesos. Su aplicación es común en sistemas de control de temperatura, velocidad, presión, flujo y nivel, donde se requiere mantener una variable de proceso cercana a un valor deseado denominado \textit{setpoint}. :contentReference[oaicite:0]{index=0}

El principio de funcionamiento de un controlador PID consiste en medir continuamente la diferencia entre la variable de proceso y el punto de referencia. Esta diferencia, conocida como error, es procesada mediante tres acciones de control: proporcional, integral y derivativa. La suma de estas tres acciones genera la señal de control aplicada al actuador del sistema. :contentReference[oaicite:1]{index=1}

\subsection{Sistema de Control en Lazo Cerrado}

Los controladores PID trabajan en sistemas de control de lazo cerrado, en los cuales existe retroalimentación continua entre la salida y la entrada del sistema. En este tipo de sistemas, un sensor mide la variable de proceso y envía dicha información al controlador, el cual compara el valor medido con el valor deseado para calcular el error y corregirlo. :contentReference[oaicite:2]{index=2}

Entre los elementos principales de un sistema PID se encuentran:

\begin{itemize}
    \item \textbf{Variable de proceso (PV):} magnitud que se desea controlar.
    \item \textbf{Setpoint (SP):} valor deseado para la variable de proceso.
    \item \textbf{Error:} diferencia entre el setpoint y la variable de proceso.
    \item \textbf{Actuador:} elemento encargado de modificar el sistema físico.
    \item \textbf{Sensor:} dispositivo que mide la variable de proceso.
\end{itemize}

El desempeño de un sistema de control suele evaluarse mediante parámetros como tiempo de subida, sobreimpulso, tiempo de establecimiento y error en estado estacionario. :contentReference[oaicite:3]{index=3}

\subsection{Acción Proporcional}

La acción proporcional genera una respuesta directamente proporcional al error presente en el sistema. Su expresión matemática es:

\[
u_P(t)=K_p e(t)
\]

donde:

\begin{itemize}
    \item $u_P(t)$ es la salida proporcional,
    \item $K_p$ es la ganancia proporcional,
    \item $e(t)$ es el error del sistema.
\end{itemize}

Un aumento en la ganancia proporcional mejora la rapidez de respuesta del sistema; sin embargo, valores excesivos pueden provocar oscilaciones e inestabilidad. :contentReference[oaicite:4]{index=4}

\subsection{Acción Integral}

La acción integral considera la acumulación del error a lo largo del tiempo, permitiendo eliminar el error en estado estacionario. Matemáticamente se expresa como:

\[
u_I(t)=K_i \int e(t)\,dt
\]

donde $K_i$ representa la ganancia integral.

La acción integral mejora la precisión del sistema, aunque un valor elevado puede incrementar el sobreimpulso y producir respuestas lentas o inestables. Además, puede aparecer el fenómeno conocido como \textit{integral windup}, el cual ocurre cuando la acción integral se acumula excesivamente. :contentReference[oaicite:5]{index=5}

\subsection{Acción Derivativa}

La acción derivativa predice el comportamiento futuro del error mediante su tasa de cambio. Su ecuación es:

\[
u_D(t)=K_d \frac{de(t)}{dt}
\]

donde $K_d$ es la ganancia derivativa.

Esta acción contribuye a reducir el sobreimpulso y mejorar la estabilidad del sistema. Sin embargo, es sensible al ruido presente en las señales medidas, por lo que generalmente se emplean valores pequeños de ganancia derivativa. :contentReference[oaicite:6]{index=6}

\subsection{Ecuación General del Controlador PID}

La combinación de las tres acciones da lugar a la ecuación general del controlador PID:

\[
u(t)=K_p e(t)+K_i \int e(t)\,dt+K_d \frac{de(t)}{dt}
\]

En el dominio de Laplace, la función de transferencia del controlador PID se expresa como:

\[
G(s)=K_p+\frac{K_i}{s}+K_d s
\]

Esta estructura permite obtener sistemas rápidos, estables y con mínimo error permanente. :contentReference[oaicite:7]{index=7}

\subsection{Sintonización de un Controlador PID}

La sintonización de un controlador PID consiste en ajustar adecuadamente las ganancias proporcional, integral y derivativa para lograr una respuesta óptima del sistema. Entre los métodos más utilizados se encuentra el método de Ziegler–Nichols, el cual determina los parámetros PID a partir de la ganancia crítica y el período de oscilación del sistema. :contentReference[oaicite:8]{index=8}

El ajuste correcto de los parámetros PID permite obtener:

\begin{itemize}
    \item menor tiempo de respuesta,
    \item reducción del error permanente,
    \item disminución del sobreimpulso,
    \item mayor estabilidad y robustez del sistema.
\end{itemize}

\subsection{Aplicaciones de los Controladores PID}

Los controladores PID son ampliamente utilizados en la automatización industrial y en sistemas embebidos debido a su facilidad de implementación y eficiencia. Sus aplicaciones incluyen:

\begin{itemize}
    \item control de temperatura en hornos y cámaras térmicas,
    \item control de velocidad en motores eléctricos,
    \item regulación de presión y caudal,
    \item sistemas robóticos,
    \item control de procesos químicos e industriales.
\end{itemize}

Gracias a su versatilidad, los controladores PID continúan siendo una de las herramientas más importantes dentro de la ingeniería de control. :contentReference[oaicite:9]{index=9}



% =============================================================================

% -----------------------------------------------------------------------------

\section{Consideraciones tecnicas del codigo PID en Controllino}

\subsection{Plataforma y estructura general}
La implementacion se realizo sobre una placa \textbf{Controllino Mega}.
El control de velocidad del motor DC se organiza en dos niveles:
\begin{itemize}
  \item \textbf{ISR de Timer1}: adquiere RPM y ejecuta el PID discreto cada periodo de muestreo.
  \item \textbf{Bucle principal}: gestiona HMI, setpoint, seguridad (emergencia/continuar) y visualizacion.
\end{itemize}

\subsection{Definicion del tiempo de muestreo y ticks de Timer1}
En el codigo se define:
\begin{itemize}
  \item Frecuencia de muestreo: $f_s = 20\,\mathrm{Hz}$
  \item Periodo de muestreo: $T = 1/f_s = 0.05\,\mathrm{s}$
\end{itemize}

El Timer1 se configura con preescalador 256.
Para un reloj base de $16\,\mathrm{MHz}$:
\[
 f_{\text{timer}} = \frac{16\times10^6}{256} = 62500\,\mathrm{ticks/s}
\]

El valor de comparacion se fija con:
\[
\text{OCR1A} = \frac{62500}{f_s} = \frac{62500}{20} = 3125
\]

Por tanto, \textbf{no son 6000 ticks} en esta configuracion, sino aproximadamente
\textbf{3125 ticks por muestra} (20 Hz).

\subsection{Medicion de velocidad (RPM)}
En cada interrupcion de muestreo se usan los pulsos contados por el encoder:
\[
\mathrm{RPM}(k)=\frac{N_p(k)\cdot 60\cdot f_s}{\mathrm{PPR}}
\]
donde:
\begin{itemize}
  \item $N_p(k)$: pulsos acumulados en la ventana de muestreo
  \item $\mathrm{PPR}=36$: pulsos por revolucion
  \item $f_s=20\,\mathrm{Hz}$
\end{itemize}

\subsection{PID discreto implementado en el codigo}
La forma usada en el firmware es la ecuacion incremental por diferencias:
\[
 u(k)=u(k-1)+K_p\left[(e(k)-e(k-1)) + \frac{T}{T_i}e(k) + \frac{T_d}{T}(e(k)-2e(k-1)+e(k-2))\right]
\]
con:
\[
 e(k)=r(k)-y(k)
\]

\textbf{Ventaja de esta forma}: evita recalcular la salida absoluta completa y usa memoria
minima del estado previo ($u(k-1), e(k-1), e(k-2)$).

\subsection{Relacion con la forma por ganancias (Ki, Kd)}
Si se desea expresar el controlador en terminos de ganancias:
\[
K_i=\frac{K_p}{T_i},\qquad K_d=K_p\,T_d
\]
la misma ecuacion puede escribirse como:
\[
 u(k)=u(k-1)+K_p(e(k)-e(k-1)) + K_i\,T\,e(k) + \frac{K_d}{T}(e(k)-2e(k-1)+e(k-2))
\]

Tambien puede representarse en forma de coeficientes:
\[
 u(k)=u(k-1)+a_0 e(k)+a_1 e(k-1)+a_2 e(k-2)
\]
con:
\[
 a_0=K_p+K_iT+\frac{K_d}{T},\quad
 a_1=-K_p-2\frac{K_d}{T},\quad
 a_2=\frac{K_d}{T}
\]

\subsection{Saturacion y escalado de la accion de control}
La salida del PID se limita y luego se escala a PWM:
\begin{itemize}
  \item Saturacion intermedia: $u\in[0,5]$
  \item Mapeo posterior: $u\in[0,255]$ para \texttt{analogWrite}
\end{itemize}

Este escalado permite compatibilizar la magnitud calculada del control con el rango real de
accionamiento del pin PWM.

\subsection{Sintonizacion y ajuste fino realizados}
En las pruebas se consideraron tres etapas:
\begin{enumerate}
  \item Valores iniciales teoricos/calculados.
  \item Ajustes intermedios obtenidos con apoyo de herramientas (p.ej. Matlab/LabVIEW).
  \item Ajuste fino empirico en planta real para reducir oscilacion y mejorar seguimiento.
\end{enumerate}

Parametros finales usados en el codigo funcional:
\[
K_p=0.03,\qquad T_i=0.08\,\mathrm{s},\qquad T_d=0.00\,\mathrm{s}
\]

Interpretacion practica del ajuste fino:
\begin{itemize}
  \item Se redujo el termino derivativo a cero para evitar amplificacion de ruido.
  \item Se uso una accion proporcional moderada para disminuir sobreimpulso.
  \item Se mantuvo la accion integral para eliminar error en regimen permanente.
\end{itemize}

\subsection{Seguridad y operacion}
El control incluye logica de seguridad por botones fisicos:
\begin{itemize}
  \item \textbf{Emergencia}: fuerza PWM a cero y detiene el motor.
  \item \textbf{Continuar}: restablece la aplicacion del control PID.
\end{itemize}

Esto asegura una parada inmediata ante condicion de riesgo y recuperacion controlada del proceso.





% =============================================================================
\section{Respuesta a las Preguntas de Analisis}
% =============================================================================

\subsection*{P1. ¿Cómo se seleccionaron los parámetros $Kp$, $Ki$ y $Kd$? Justificar el método de sintonización utilizado y explicar el efecto de cada ganancia en la respuesta del sistema.}
\section{P1 — Selección de los parámetros $K_p$, $K_i$ y $K_d$}

\subsection{Método de sintonización utilizado}

La selección de los parámetros del controlador PID se realizó en tres etapas
sucesivas:

\begin{enumerate}
  \item \textbf{Valores iniciales teóricos.}
        A partir de la función de transferencia del motor obtenida en el
        laboratorio previo (mediante adquisición con LabVIEW o Matlab/PLX-DAQ),
        se estimó un punto de partida usando reglas analíticas de diseño en
        tiempo discreto.

  \item \textbf{Ajuste asistido por herramientas.}
        Con la función de transferencia identificada se realizó una simulación
        en Matlab (o Simulink) para verificar el comportamiento en lazo cerrado
        y refinar las ganancias antes de cargar el código al Controllino.

  \item \textbf{Ajuste fino empírico en planta real.}
        Los parámetros finales se ajustaron directamente en la planta
        observando la respuesta en la HMI (serie de RPM vs. referencia) y
        minimizando sobreimpulso, oscilaciones y error en régimen permanente.
\end{enumerate}

\subsection{Parámetros finales implementados}

Los valores finales cargados en el firmware Stone son:
\[
  K_p = 0.03, \qquad T_i = 0.08\,\text{s}, \qquad T_d = 0\,\text{s}
\]

Lo que equivale a un \textbf{controlador PI}. La Tabla~\ref{tab:pid_params}
resume los parámetros para ambas versiones de HMI.

\begin{table}[H]
  \centering
  \caption{Parámetros PID implementados por plataforma HMI}
  \label{tab:pid_params}
  \begin{tabular}{lccc}
    \toprule
    \textbf{Plataforma} & $K_p$ & $K_i$ & $K_d$ \\
    \midrule
    Stone (Ti/Td)   & 0.03 & $K_p/T_i = 0.375$ & 0 \\
    Coolmay (Kp/Ki/Kd) & 0.30 & 0.08 & 0.01 \\
    \bottomrule
  \end{tabular}
\end{table}

\subsection{Efecto de cada ganancia sobre la respuesta}

\begin{itemize}
  \item \textbf{$K_p$ (acción proporcional):} reduce el error proporcional al
        error actual. Un valor elevado disminuye el tiempo de subida pero
        incrementa el sobreimpulso y puede desestabilizar el sistema. En la
        implementación, $K_p$ es moderado para mantener la respuesta suave.

  \item \textbf{$K_i$ (acción integral):} acumula el error en el tiempo y lo
        elimina en régimen permanente. Es la razón por la que el sistema
        alcanza exactamente el setpoint de RPM sin error estacionario. Un $K_i$
        excesivo provoca sobreimpulso prolongado e \textit{integral windup}.

  \item \textbf{$K_d$ (acción derivativa):} anticipa la trayectoria del error
        y frena la acción antes de superar el setpoint. En el ajuste final se
        fijó $T_d = 0$ (versión Stone) para evitar la amplificación de ruido
        en la señal de RPM procedente del encoder.
\end{itemize}

La ecuación incremental discretizada implementada en la ISR del Timer1 es:
\begin{equation}
  u(k) = u(k-1) + a_0\,e(k) + a_1\,e(k-1) + a_2\,e(k-2)
  \label{eq:pid_incr}
\end{equation}
con coeficientes:
\begin{align}
  a_0 &= K_p + K_i T + \frac{K_d}{T} \\
  a_1 &= -K_p - \frac{2K_d}{T} \\
  a_2 &= \frac{K_d}{T}
\end{align}
donde $T$ es el período de muestreo.



\subsection*{P2. ¿Cuál es el tiempo de muestreo seleccionado para el controlador PID? Justificar la elección del tiempo de muestreo y explicar cómo se calculan las interrupciones para mantener ese tiempo de muestreo.}

\subsection{Tiempo de muestreo seleccionado}

Para la versión Stone se seleccionó:
\[
  f_s = 20\,\text{Hz} \qquad \Longrightarrow \qquad T = \frac{1}{f_s} = 0.05\,\text{s}
\]

\textbf{Justificación:} La dinámica de velocidad de un motor DC tiene
constantes de tiempo del orden de decenas de milisegundos. Un período de
muestreo de $T = 50\,\text{ms}$ garantiza que:
\begin{itemize}
  \item Se muestrea al menos $5{-}10$ veces por constante de tiempo del motor,
        cumpliendo el criterio práctico de diseño digital.
  \item El conteo de pulsos del encoder acumula un número suficiente de pulsos
        por ventana para obtener una estimación de RPM con resolución adecuada.
  \item La carga computacional en la ISR es baja, permitiendo al bucle
        principal atender la comunicación con el HMI.
\end{itemize}

\subsection{Cálculo de las interrupciones del Timer1}

El Timer1 del Controllino (ATmega2560, reloj a $16\,\text{MHz}$) se configura
en modo CTC con preescalador 256:

\[
  f_{\text{timer}} = \frac{f_{\text{clk}}}{\text{preescalador}}
  = \frac{16 \times 10^6}{256} = 62\,500\,\text{ticks/s}
\]

El valor de comparación \texttt{OCR1A} se calcula de modo que la interrupción
ocurra exactamente cada $T$ segundos:

\[
  \text{OCR1A} = \frac{f_{\text{timer}}}{f_s} = \frac{62\,500}{20} = 3\,125
\]

En código esto se expresa como:
\begin{lstlisting}[language=C++, caption={Configuración del Timer1 a 20\,Hz (Stone PID)}]
noInterrupts();
TCCR1A = 0;
TCCR1B = 0b00000100;   // prescaler 256
TIMSK1 |= B00000010;   // interrupciones por comparacion
OCR1A  = 62500 / fs;   // 62500 / 20 = 3125 ticks
interrupts();
\end{lstlisting}

Cada vez que \texttt{TCNT1} alcanza \texttt{OCR1A}, se lanza
\texttt{ISR(TIMER1\_COMPA\_vect)}, que calcula las RPM y ejecuta el PID:

\[
  \text{RPM}(k) = \frac{N_p(k)\cdot 60 \cdot f_s}{\text{PPR}}
  = \frac{N_p(k)\cdot 60 \cdot 20}{36} = N_p(k)\cdot 33.33
\]

donde $N_p(k)$ son los pulsos contados por el encoder en la ventana de
muestreo y $\text{PPR} = 36$ es la resolución del encoder.



\subsection*{P3. ¿El controlador actúa de forma efectiva en la toda la zona lineal del motor? Justificar la respuesta con base en la función de transferencia obtenida y la respuesta del sistema.}


\subsection{Definición de la zona lineal}

Un motor DC real presenta saturación en el accionamiento (valores de PWM muy
bajos no producen movimiento por la fricción estática) y saturación a plena
velocidad. Entre ambos extremos existe una \textbf{zona aproximadamente lineal}
donde la velocidad responde de forma proporcional al ciclo de trabajo (duty
cycle) aplicado.

La función de transferencia de primer orden identificada experimentalmente
adopta la forma:
\[
  G(s) = \frac{K_m}{\tau_m s + 1}
\]
donde $K_m$ es la ganancia estática (RPM/unidad de PWM) y $\tau_m$ es la
constante de tiempo mecánica del motor.

\subsection{Comportamiento observado}

La respuesta del sistema en lazo cerrado presentó:
\begin{itemize}
  \item \textbf{Seguimiento de referencia:} el controlador PI ($T_d = 0$)
        alcanza el setpoint de RPM con error estacionario nulo gracias a la
        acción integral, incluso ante perturbaciones de carga moderadas.
  \item \textbf{Sobreimpulso reducido:} el valor moderado de $K_p = 0.03$
        (versión Stone) mantuvo el sobreimpulso dentro de un rango aceptable
        ($<10\%$) en los puntos de operación dentro de la zona lineal.
  \item \textbf{Estabilidad:} la respuesta fue estable y convergente para
        referencias dentro del rango operativo del motor (aprox.\
        200–3500\,RPM), que corresponde a la zona lineal identificada.
\end{itemize}

\subsection{Limitaciones fuera de la zona lineal}

En los extremos de operación se observó:
\begin{itemize}
  \item \textbf{Velocidades muy bajas} ($<\!\!200\,\text{RPM}$): la fricción
        estática introduce una no linealidad que el controlador no puede
        compensar completamente. El motor puede detenerse aunque la señal de
        control $u(k) > 0$.
  \item \textbf{Velocidades muy altas} (próximas al máximo): la acción de
        control satura en $u = 255$ (límite del PWM de 8~bits), degradando el
        comportamiento en lazo cerrado ya que el integrador continúa
        acumulando (\textit{windup}).
  \item \textbf{Perturbaciones bruscas de carga:} fuera del punto de
        linealización, la ganancia del proceso cambia y los parámetros del PID
        sintonizados para el punto nominal pueden resultar subóptimos.
\end{itemize}


% =============================================================================
\section{Conclusiones}
% =============================================================================


\begin{itemize}
    \item Se logró implementar correctamente un sistema de control y monitoreo de velocidad para un motor DC utilizando el Controlino Mega y una interfaz HMI.

    \item El uso de interrupciones externas y temporizadores permitió realizar una adquisición precisa de los pulsos del encoder y un cálculo confiable de la velocidad del motor.

    \item La señal PWM demostró ser una técnica eficiente para regular la velocidad del motor de corriente continua de manera estable y controlada.

    \item La integración de la HMI facilitó la visualización en tiempo real de las variables del proceso y permitió ajustar parámetros de control de forma intuitiva.

    \item La implementación del controlador PID mejoró la respuesta dinámica del sistema, reduciendo el error entre la referencia y la velocidad real del motor.

    \item La práctica permitió reforzar conocimientos relacionados con control digital, sistemas embebidos, adquisición de datos y automatización industrial.
\end{itemize}

% =============================================================================
\addcontentsline{toc}{section}{Referencias}
% =============================================================================

\subsection{Referencias}

\begin{itemize}
    \item National Instruments. ``Explicación sobre el controlador PID y la teoría''. Disponible en: \\
    [NI PID Theory Explained](https://www.ni.com/es/shop/labview/pid-theory-explained.html?utm_source=chatgpt.com)
    
    \item LibreTexts Español. ``Controladores PI, PD y PID''. Disponible en: \\
    [LibreTexts Controladores PID](https://espanol.libretexts.org/Bookshelves/Ingenieria/Ingenieria_Industrial_y_de_Sistemas/Libro%3A_Introducci%C3%B3n_a_los_Sistemas_de_Control_%28Iqbal%29/03%3A_Modelos_de_sistemas_de_control_de_retroalimentaci%C3%B3n/3.3%3A_Controladores_PI%2C_PD_y_PID?utm_source=chatgpt.com)
\end{itemize}

% =============================================================================
%  ANEXOS
% =============================================================================

\newpage
\appendix

\section{Código Fuente --- Base: Control de Motor DC}

El siguiente listado corresponde al archivo \texttt{Practica3\_Stone\_Base.ino},
que implementa el control de un motor DC mediante un slider de Duty Cycle. Además, implementa una gráfica donde se puede observar las RPM del motor y su duty cycle ingresado.


\lstinputlisting[style=ArduinoStyle]{Practica3_Stone.ino}


% -----------------------------------------------------------------------------

\section{Código Fuente --- Reto: Control de Motor DC con PID}

El siguiente listado corresponde al archivo \texttt{Practica3\_Stone\_PID.ino},
que añade el controlador PID al motor DC.

\lstinputlisting[style=ArduinoStyle] {Practica3_Stone_Reto.ino}


\end{document}
